\documentclass[a4paper,11pt]{article}
\pdfoutput=1 % if your are submitting a pdflatex (i.e. if you have
             % images in pdf, png or jpg format)
\usepackage{jinstpub} % for details on the use of the package, please
                     % see the JINST-author-manual

\title{REST-for-Physics, a ROOT-based framework for event oriented data analysis and combined Monte Carlo response.}

\collaboration{\includegraphics[height=24mm]{RESTlogoFull.png}}

%% %simple case: 2 authors, same institution
%% \author{A. Uthor}
%% \author{and A. Nother Author}
%% \affiliation{Institution,\\Address, Country}

% more complex case: 4 authors, 3 institutions, 2 footnotes
\author[a,b,1]{F. Irst,\note{Corresponding author.}}
\author[c]{S. Econd,}
\author[a,2]{T. Hird\note{Also at Some University.}}
\author[c,2]{and Fourth}

% The "\note" macro will give a warning: "Ignoring empty anchor..."
% you can safely ignore it.

\affiliation[a]{One University,\\some-street, Country}
\affiliation[b]{Another University,\\different-address, Country}
\affiliation[c]{A School for Advanced Studies,\\some-location, Country}

% e-mail addresses: only for the forresponding author
\emailAdd{javier.galan@unizar.es}
\emailAdd{bur\_ning@sjtu.edu.cn}

\abstract{Abstract...}


\keywords{Software architectures (event data models, frameworks and databases);Simulation methods and programs;Data processing methods;}

%\arxivnumber{1234.56789} % only if you have one


% \collaboration{\includegraphics[height=17mm]{example-image}\\[6pt]
%   XXX collaboration}
% or
%\collaboration[c]{on behalf of XXX collaboration}

% if you write for a special issue this may be useful
%\proceeding{N$^{\text{th}}$ Workshop on X\\
%  when\\
%  where}

\begin{document}
\maketitle{}
\flushbottom
%%%%% I create sections, but it is just orientative. Some sections may disappear and get integrated in other sections. It is just to do not forget to talk about something.

\input{introduction}
\input{description}

%section{Code structure and data traceability.}
%\input{03_structure}

% Libries and packages could be merged into one?
\input{libraries}

\input{summary}
\acknowledgments
We acknowledge support from the the European Research Council (ERC) under the European Union’s Horizon 2020 research and innovation programme, grant agreement ERC-2017-AdG788781 (IAXO+), and from the Spanish Agencia Estatal de Investigacion under grant FPA2016-76978-C3-1-P.

%\begin{figure}[htb!]
  \centering
  \raisebox{-0.5\height}{\includegraphics[width=0.99\linewidth]{images/institution_logos.png}}
%\end{figure}

%\section*{References}
\bibliography{paper}
\bibliographystyle{unsrt}

% We suggest to always provide author, title and journal data:
% in short all the informations that clearly identify a document.

%\begin{thebibliography}{99}

%\bibitem{a}
%Author, \emph{Title}, \emph{J. Abbrev.} {\bf vol} (year) pg.

%\bibitem{b}
%Author, \emph{Title},
%arxiv:1234.5678.

%\bibitem{c}
%Author, \emph{Title},
%Publisher (year).


% Please avoid comments such as "For a review'', "For some examples",
% "and references therein" or move them in the text. In general,
% please leave only references in the bibliography and move all
% accessory text in footnotes.

% Also, please have only one work for each \bibitem.



\end{document}